\documentclass[12pt]{article}
\usepackage{amsmath, amssymb, amsthm}
\usepackage[margin=1in]{geometry}

\title{Problem 265: Binary Circles}
\author{Project Euler Solution}
\date{}

\begin{document}
\maketitle

\section{Problem Statement}
A binary circle of order $N$ is a circular arrangement of $2^N$ binary digits such that all $N$-bit clockwise subsequences are distinct. For $N=5$ (32 bits, 32 distinct 5-bit subsequences), how many such arrangements exist?

We fix the first 5 bits as \texttt{00000} to remove rotational equivalence.

\section{Solution}

\subsection{De Bruijn Graph}
Construct the de Bruijn graph $G(2, 4)$: vertices are all 4-bit strings, edges represent 5-bit strings. Each vertex has in-degree~2 and out-degree~2.

A binary circle of order 5 corresponds to an Eulerian circuit in this graph starting from vertex \texttt{0000}.

\subsection{Enumeration}
We enumerate all Eulerian circuits via DFS with backtracking:
\begin{enumerate}
\item Start at vertex \texttt{0000}, edge \texttt{00000} (bit 0).
\item At each step, choose to append 0 or 1.
\item Track which edges (5-bit strings) have been used.
\item A complete circuit uses all 32 edges and returns to \texttt{0000}.
\end{enumerate}

\subsection{Complexity}
The backtracking explores at most $O(2^{32})$ paths but heavy pruning (dead-end detection) makes it feasible.

\section{Answer}

\[
\boxed{209110240768}
\]

\end{document}
